\section{The complete ordinary automorphy-away endpoint}
\label{sec:automorphy-away-v45}

We retain the original $H'=\mathrm{HypothesesH'}$ and all its
named literature inputs.  In particular, the definite-inner-form
base-change input and polarized cuspidal attachment remain the
ones used in BCS and GUA.  The local inputs are the original
admissibility and Jacquet finite-length theorem, the geometric
lemma and uniqueness of cuspidal support, the generic unlinked
segment classification, smooth contragredient duality and
Gelfand--Kazhdan, and the full $\mathrm{IsRecT}$ laws and existence
statement.  We use the already matched generic-unitary theorem
BCS8, not a new exponent estimate.  Global strong multiplicity one,
Chebotarev, local/global reciprocity, and the basic central-character
and compact norm-one idele-class facts are retained foundations.
Below we prove their finite consumers, including the single global
unitarizing twist; in
particular the window conclusion is not retained as an opaque
extra assumption.

Fix the original $\ell,L,\Omega,s,n,V,Q,c,r,\iota,B,\Lambda,R,K$,
where $n\ge1$, $R$ is the given unitary Hecke datum, and
$\iota:\overline{\mathbb Q}_\ell\simeq\mathbb C$ is the given
field isomorphism.  The final endpoint assumes $V.S'=s.S$ and
$R.\mathrm{IsJ}(B,K,J)$, with exactly the same $J$.  No
$r.\mathrm{V0ok}$ hypothesis is added.  Write $\varepsilon$ for
the cyclotomic character.  At a finite place $w$, set
$q_w=\#(\mathcal O_L/w)>1$, and choose the source's uniformizer
$\varpi_w=\mathrm{unif}(L_w)$.  The absolute value satisfies
$|\varpi_w|_w=q_w^{-1}$.

Here is the exact local vocabulary.  A cuspidal support multiset
$S$ is \emph{toral} when all its cuspidal entries have rank one.
Its smooth rank-one characters form $\mathrm{chars}(S)$, with
the original multiplicities retained in $S$.  Put
\[
 \mathrm{Unl}(q,S):\quad
 \xi,\xi'\in\mathrm{chars}(S),\quad \xi/\xi'\text{ unramified}
       \ \Longrightarrow\ (\xi/\xi')(\varpi)\ne q .
 \tag{HAW.1}
\]
For an actual cusp $\pi$, $\mathrm{TU}(\pi,w)$ means
$\mathrm{TorAt}(\pi,w)$ and (HAW.1) for every one of its
support witnesses.  The literal strict window is
\[
 \begin{gathered}
 \mathrm{WinAt}(\pi,t,w):\quad \mathrm{TorAt}(\pi,w),\qquad
 q_w^{t-1/2}<|\xi(\varpi_w)|<q_w^{t+1/2},\\
 S\text{ a support at }w,\qquad\xi\in\mathrm{chars}(S).
 \end{gathered}
 \tag{HAW.2}
\]
For a Galois representation $\rho$, $\mathrm{Unl1}(\rho,w)$
requires at \emph{every} prolongation over $w$ an actual full flag
with characters $\kappa_i$ satisfying
$\kappa_i\ne\kappa_j\varepsilon$ for all $i,j$.
$\mathrm{PotUnl}(\rho,w)$ requires instead
\[
       \kappa_i^m\ne(\kappa_j\varepsilon)^m
                         \quad(i,j;\ m\ge1).
 \tag{HAW.3}
\]
All flag orders, principal-series normalizations, and full
determinant polynomial laws are the original ones.

\subsection{The same automorphic part and all places away from the auxiliary prime}

% BEGIN HAW1
\paragraph{HAW1. Comparison with an arbitrary given representation.}
\label{claim:HAW1}
Let $\sigma$ be an original actual $\mathrm{IsAutRepU}$ space
over a CM extension unramified at every finite place, and let
$\rho:G_L\to\mathrm{GL}_n(\overline{\mathbb Q}_\ell)$ be
continuous and irreducible.  Suppose that, outside one finite set,
at every split place every actual local component of $\sigma$
and each of its Satake multisets gives the original
$\mathrm{FrobEig}$ formula for this $\rho$.  Then there is an
actual cusp $\pi$ with
\[
 \mathrm{RegAlg}(\pi),\quad \mathrm{DualIsConj}(\pi),\quad
 \mathrm{AttC}(\iota,\rho,\pi),
 \tag{HAW.4}
\]
Satake parameters at every $\mathrm{IsUnramU}(\sigma,w)$,
and, at every split $w$, the equivalence
\[
 \pi.\mathrm{IsLocComp}(w,\tau)
 \ \Longleftrightarrow\
 \mathrm{IsCompAlong}(\sigma,\mathrm{locU},\mathrm{cLoc},\tau).
 \tag{HAW.5}
\]
Here $\tau$ ranges over every actual smooth irreducible representation.
This is the literal same-given-representation $\mathrm{link2\_alt}$.
\begin{proof}
GUA18 supplies a continuous semisimple representation $r_\sigma$
with the full local comparison and GUA17's irreducible-cusp
implication.  Intersect the given cofinite set with the places
away from $\ell$ and those where $\sigma$ is unramified.
An absolute residue-degree-one place in this set is split:
its relative residue degree is one, and unramifiedness in the
quadratic extension makes the degree formula force two places
above the real place.  At it choose an actual local component
by LTP19.  It is spherical by LTP20, so it has a Satake multiset.
UPA7 and the given hypothesis identify the characteristic
polynomials of $r_\sigma$ and $\rho$ at all geometric Frobenius
lifts there.

CAP6's degree-one Chebotarev density extends this equality to
every element of $G_L$: the polynomial coefficients are
continuous over $\overline{\mathbb Q}_\ell$, a Hausdorff field.
The field isomorphism $\iota$ is used only after those
equalities; no complex continuity of $\iota$ is asserted.
The two semisimple representations are isomorphic.  The finite
algebra argument also applies later without any topological
hypothesis: take the image of the group algebra in
$\operatorname{End}(W\oplus W')$.  Its radical kills the
semisimple faithful module, hence is zero.  It is a product of
matrix algebras over the algebraically closed coefficient field.
Equality of traces on group elements gives equality on this
algebra; evaluating its central block idempotents determines
the multiplicity of every simple module.  The representations
are therefore isomorphic.  Characteristic zero makes division
by each positive simple dimension legitimate.

Thus $r_\sigma$ is irreducible.  GUA17 supplies $\pi$ and
(HAW.4), transferring the attachment across this isomorphism.
It also supplies the Satake assertion at every unramified place
and the forward implication in (HAW.5).  For its reverse,
choose one actual local component $\tau'$ of $\pi$ at $w$
by BCS4.
Both $\tau'$ and the supplied $\tau$ are components of the same
$\sigma$.  Uniqueness of the split actual component gives
$\tau\simeq\tau'$, so $\tau$ is a component of $\pi$.
This keeps the supplied $\rho$ and proves all the local
quantifiers, not just a cofinite local assertion.
\end{proof}
% END HAW1

% BEGIN HAW2
\paragraph{HAW2. The literal PartLinkAlt and the partwise ordinary cusp predicate.}
\label{claim:HAW2}
Under the original $H'$, every nonzero $YOf$ in the
$\mathrm{PartLinkAlt}$ contract has an original
$Y:\mathrm{PartAltAlt}$ with $Y.Y=YOf$ such that irreducibility
of $Y.r\circ V.\mathrm{pr}$ gives (HAW.4), the entire
unramified Satake assertion, and the local comparison in
(HAW.5).  Moreover the original $\mathrm{AutUOAlt}$ holds:
the actual classical spaces are spanned by parts whose
irreducible pullbacks have such a cusp with
$\pi.\mathrm{IsLocComp}(s.w(v),Y.\sigma_l(v))$ for every
$v\mid\ell$.
\begin{proof}
Use UPA16 on the supplied type, weight, coefficient basis,
algebraic representation, and the given actual $\sigma$.
It constructs the original $Y$ with exactly $Y.Y=YOf$, a
continuous semisimple frame representation, and full split
local comparison.  Pull it back by the original continuous
$V.\mathrm{pr}$.  UPA7 converts that comparison to the
cofinite Frobenius premise of HAW1.  If the pullback is
irreducible, HAW1 applies without reselecting $Y$.
This proves all fields of $\mathrm{PartLinkAlt}$.

For $\mathrm{AutUOAlt}$ use ADP12's decomposition of the
actual vector into these $YOf$ spaces.  The given integral
type has open kernel, as in ITA12--ITA13, so that decomposition
applies with the original finite type.  Zero parts contribute
nothing.  The nonzero parts have the preceding predicate.
For each chosen $v\mid\ell$, the part's $Y.\sigma_l(v)$
is an actual component of $\sigma$.  Formula (HAW.5) says it
is a component of the same $\pi$.  Consequently the spanning
statement has the full predicate of $\mathrm{AutUOAlt}$,
including its dependence on the original part $Y$.
\end{proof}
% END HAW2

% BEGIN HAW3
\paragraph{HAW3. The plain-support dictionary and its original type supplier.}
\label{claim:HAW3}
For every positive-rank actual cusp, all its local components
are $\mathrm{PlainTor}$ at $w$ if and only if
$\mathrm{TU}(\pi,w)$.  A spherical component of a cusp has
this property.  In addition, the full original
$\mathrm{TypePlain}$ implication holds: if the parameter of
the algebraically normalized local component has zero
monodromy and the specified inertia characteristic polynomial
of restrictions of global characters, then that local
component is $\mathrm{PlainTor}$.
\begin{proof}
Write $\nu=|\cdot|_w$.  Two smooth characters have an
unramified quotient determined by its value on $\varpi_w$,
because every element of $L_w^\times$ is a uniformizer power
times a unit.  Since $\nu(\varpi_w)=q_w^{-1}$,
\[
 [\xi]=[\xi']\nu
 \quad\Longleftrightarrow\quad
 \xi'/\xi\text{ is unramified and }
                       (\xi'/\xi)(\varpi_w)=q_w.
 \tag{HAW.6}
\]
Uniqueness of cuspidal support identifies every support witness
and every isomorphic actual component.  Its rank-one entries
are characters.  Thus (HAW.6), with the ordered pair reversed
when necessary, is exactly the equivalence between the source's
no-$\nu$-linked-pair condition and (HAW.1).  This proves the
first assertion with all support witnesses, not just one.

Here is the finite converse from the supplied inertia type.
Apply the retained Langlands multisegment classification and
the $\mathrm{IsRecT}$ sum law.  For any segment beginning in
a supercuspidal $\delta$, the characteristic polynomial of
its irreducible Weil parameter divides that of the total
parameter on every Weil element.  At a commutator
$[\sigma,t]$, $\sigma$ Weil and $t$ inertia, the stipulated
global characters all equal one.  This commutator lies in
inertia, so the type hypothesis makes the total polynomial
$(X-1)^n$.  The constituent parameter is semisimple on
inertia, and a semisimple matrix with characteristic polynomial
dividing $(X-1)^n$ is the identity.  Hence every inertia
matrix of this irreducible Weil constituent commutes with the
entire Weil image.  By the elementary Schur argument it is
scalar: an eigenspace of a commuting matrix is a nonzero
invariant subspace.  Choose an eigenline of Frobenius.  Inertia
acts on it by scalars, and inertia together with Frobenius
generates the Weil group, so irreducibility forces dimension
one.  The correspondence's rank equality then forces every
cuspidal support entry to have rank one.  This is the full
abelian-converse consumer needed here; it uses the original
global characters, not arbitrary inertia characters.

The actual local component of a cusp is generic (BCS5).
The generic classification writes its algebraic normalization
as the irreducible product of mutually unlinked segments.
Zero monodromy in the parameter forces every segment to have
length one: a segment of length greater than one contributes
its nonzero standard nilpotent block, and a direct sum cannot
cancel that block.  Two distinct singleton characters whose
ratio is $\nu$ are linked segments, a contradiction.  A
character cannot equal its own $\nu$ twist, since their
uniformizer values differ by $q_w^{-1}\ne1$.
Thus the support is plain.  The common algebraic normalization
twist changes neither the ranks nor character ratios, proving
the stated $\mathrm{TypePlain}$ implication for the original
component and every permitted correspondence.

Finally a spherical generic irreducible has unramified
rank-one cuspidal support and zero-monodromy parameter, by
GUA12 and its spherical classification argument.  The preceding
generic singleton argument gives plain support.  Equivalently,
GUA12's character support combined with the unlinked generic
classification excludes the same adjacent pair.  This proves
the original $\mathrm{SphUnl}$ assertion.
\end{proof}
% END HAW3

% BEGIN HAW4
\paragraph{HAW4. ConjTor, including ramified character ratios.}
\label{claim:HAW4}
If $\mathrm{DualIsConj}(\pi)$ and $w,w'$ lie over the same
place of the maximal real subfield, then
$\mathrm{TU}(\pi,w)$ implies $\mathrm{TU}(\pi,w')$.
\begin{proof}
If $w=w'$ there is nothing to prove.  Otherwise $w'=cw$.
The original duality and conjugation laws identify an actual
component at $w'$ with the inverse-transpose transport of an
actual component at $w$ through the continuous field
isomorphism $c:L_w\simeq L_{w'}$.  Cuspidal support under
this operation is obtained by contragredients and field
transport.  To justify this exact support rule, dualize a
subquotient of normalized parabolic induction.  Smooth
contragredient is exact on admissible modules: a functional
on a submodule extends linearly and then is averaged over
an open compact subgroup fixing it.  Its average still
extends the functional.  The invariant induction pairing
identifies the dual induction with induction of the dual
factors.  Field transport preserves that induction and its
cuspidal factors.  Uniqueness of support proves the rule.

The field isomorphism carries the valuation ring, its maximal
ideal, and the unit group to their counterparts.  One way to
check the first assertion without a chosen valuation scale is
that elements of the valuation ring have bounded nonnegative
powers; this property and its inverse are preserved by a
continuous field isomorphism of local fields.  The units and
nonunits then distinguish the maximal ideal.  Thus $c$ carries
a uniformizer to a uniformizer times a unit, and the residue
cardinality remains $q_w$.

Rank-one support characters transform as
$\xi\mapsto\xi^{-1}\circ c^{-1}$.  Rank one is preserved.
For a proposed forbidden pair at $w'$, take the corresponding
pair at $w$ in the reverse order.  Their quotient is the
transport of the proposed quotient; it is unramified exactly
when that quotient is unramified.  An unramified quotient
ignores the unit multiplying the transported uniformizer, so
its value is unchanged.  A value $q_{w'}=q_w$ would contradict
(HAW.1) at $w$.  This proves all the support quantifiers of
the literal $\mathrm{ConjTor}$, with no unramifiedness
assumption on the individual characters.
\end{proof}
% END HAW4

% BEGIN HAW5
\paragraph{HAW5. The complete AutAway supplier on the original spaces.}
\label{claim:HAW5}
If $V.S'=s.S$, the original $\mathrm{hset}(Q,\tau_0)$ is
an $\mathrm{AutomorphicHeckeSet}$ for the predicate
\[
 \rho\longmapsto
 [\mathrm{IsIrr}(\rho\circ V.\mathrm{pr})
       \Longrightarrow\mathrm{AwayOk}(s,\iota,
                                      \rho\circ V.\mathrm{pr})].
 \tag{HAW.7}
\]
Here $\mathrm{AwayOk}$ includes an actual regular algebraic
dual-is-conjugate attached cusp and $\mathrm{TU}(\pi,w)$
at \emph{every} $w\nmid\ell$ whose underlying real place
is not $s.v_0$.
\begin{proof}
As in HAW2, decompose every actual admissible classical vector
by ADP12, discard zero $YOf$ spaces, and apply the literal
$\mathrm{PartLinkAlt}$ to each nonzero part.  Fix a part with
irreducible pullback and its cusp $\pi$.  If the underlying
real place of $w$ is outside $V.S'$, LTP12 makes the same
automorphic space unramified there.  HAW2 supplies Satake
parameters for $\pi$ at $w$, and HAW3 gives
$\mathrm{TU}(\pi,w)$.

Otherwise $V.S'=s.S$, together with $w\nmid\ell$ and
$\underline w\ne s.v_0$, puts $\underline w$ in $s.St$.
The chosen $w_0=s.w(\underline w)$ is split.  Choose actual
local components of $\pi$ at all chosen places over $s.St$.
To prove a statement for an arbitrary component $\tau$ at
$w_0$, replace just the $w_0$ member of this family by $\tau$.
Every member is still an actual component of $\pi$, hence of
the original unitary space by HAW2.  LTP17--LTP18 and the
given $\mathrm{TypeForAlt}$ apply to this entire updated
family and every permitted $\mathrm{IsRecT}$ map.  They give
the exact zero-monodromy inertia type whose characters are
the original global Teichmuller characters restricted along
$V.\mathrm{pr}$.  HAW3 makes this arbitrary $\tau$ plain.
Its support dictionary gives $\mathrm{TU}(\pi,w_0)$; HAW4
transports that assertion to $w$.

Thus the same part has (HAW.7).  The original vector lies in
the sum of precisely these parts, which is the full
$\mathrm{AutomorphicHeckeSet}$ assertion.  Neither the
exceptional auxiliary place nor the given type was changed.
\end{proof}
% END HAW5

% BEGIN HAW6
\paragraph{HAW6. One window center for every toral unlinked place.}
\label{claim:HAW6}
For every positive-rank actual cusp $\pi$ over a totally
complex number field, there is one $t\in\mathbb R$ such that
\[
       \mathrm{TU}(\pi,w)\ \Longrightarrow\
                  \mathrm{WinAt}(\pi,t,w)
                        \quad\text{for every finite }w.
 \tag{HAW.8}
\]
This is the full original $\mathrm{WinOf}$, with its strict
half-width and its choice of $t$ before $w$.
\begin{proof}
First let $\sigma$ be a generic unitarisable local irreducible
whose support is toral and unlinked.  By BCS8 it is the actual
normalized product of unitary square-integrable segment
representations $\delta_a\nu^{s_a}$, $|s_a|<1/2$.
Support uniqueness identifies its support with the union of
the corresponding segment supports.  Every such segment has
length one.  Otherwise its first two rank-one entries would
be $\xi$ and $\xi\nu$; their reversed quotient is unramified
with uniformizer value $q$, violating (HAW.1).  Its unitary
length-one factor is therefore a unitary rank-one character
$\zeta_a$.  Every support character has the form
$\zeta_a\nu^{s_a}$, and
\[
             |\zeta_a(\varpi)q^{-s_a}|=q^{-s_a},
                     \qquad q^{-1/2}<q^{-s_a}<q^{1/2}.
 \tag{HAW.9}
\]
Repeated characters and all positive local ranks are included.

Here is the finite global unitarization step.  Let $\zeta$
be the original continuous central Hecke character of the
positive-rank cusp $\pi$.  Its positive absolute value is
trivial on the compact norm-one idele class group: bounded
powers of a point and its inverse force modulus one.
It therefore factors through the idele norm in
$\mathbb R_{>0}$.  Taking logarithms turns a continuous
homomorphism into a continuous additive function on
$\mathbb R$, hence multiplication by one real $b$.
Thus $|\zeta|=|\cdot|_{\mathbb A}^{b}$.
Choose the single character
$\chi=|\cdot|_{\mathbb A}^{-b/n}$ and the actual twisted
cusp $\pi\otimes(\chi\circ\det)$.  Its central character
has modulus $|\cdot|_{\mathbb A}^{b}\chi^n=1$.
This is the original normalized/unitary cuspidal class;
the cuspidal positive form in BCS1 makes its local
components unitarisable.  At every finite place
$|\chi_w(\varpi_w)|=q_w^{b/n}$.  Put $a=b/n$.
This proves the precise single-twist assertion of the
original global unitarization supplier and fixes $a$
before any local place is selected.

Twisting sends each support character $\xi$ to
$\xi\chi_w$ and preserves rank and every quotient of two
characters.  Hence a toral unlinked place of $\pi$ is a
toral unlinked place of $\pi\otimes\chi$.  Its actual
component is generic by BCS5 and unitarisable by BCS1, so
(HAW.9) applies.  For each support character there is
$s\in(-1/2,1/2)$ with
$|\xi(\varpi_w)|q_w^a=q_w^{-s}$.
Consequently $|\xi(\varpi_w)|=q_w^{-a-s}$, and (HAW.2)
holds with the single choice $t=-a$.  This proves (HAW.8)
for every support witness at every eligible place.
\end{proof}
% END HAW6

\subsection{Ordinary extraction at every arithmetic point}

% BEGIN HAW7
\paragraph{HAW7. A nonzero ordinary eigenvector gives the exact principal-series subquotient.}
\label{claim:HAW7}
Let $k$ be any nonarchimedean local field and let $\sigma$ be
a smooth irreducible complex representation of
$\mathrm{GL}_n(k)$, with the original Iwahori data.  Let
$0\ne v\in\mathrm{eig}(\sigma,c,\varphi)$, where
$\varphi(e_i(u))=1$ for $u\in U_c(k,c+1)$, and suppose
\[
 U_{a_j}v=b_jv,\qquad b_j\ne0\quad(1\le j\le n),\qquad
 a_j=\operatorname{diag}(\underbrace{\varpi,\ldots,\varpi}_j,1,\ldots,1).
 \tag{HAW.10}
\]
Under the original local finite-length input there are smooth
characters $\omega_i:k^\times\to\mathbb C^\times$ such that
\[
 \begin{gathered}
 \sigma\text{ is a subquotient of }\mathrm{PS}(\omega),\qquad
 \omega_i(u)=\varphi(e_i(u))^{-1},\\
 b_j=\mathrm{psEig}(\omega,j)
       =q^{-j(j-1)/2}\prod_{i<j}\omega_i(\varpi).
 \end{gathered}
 \tag{HAW.11}
\]
The normalization is exactly
$\mathrm{PS}(\omega)=\operatorname{Ind}^{\rm norm}_B
(\omega_i|\cdot|^{(n-1)/2})_i$.
\begin{proof}
For $n=0$ the group and its irreducible representation are
one-dimensional in the only possible way, and the conclusion
is immediate.  Suppose $n\ge1$.  The positive-operator
product identity of OST5 applied to
$z=a_1\cdots a_{n-1}=\operatorname{diag}
(\varpi^{n-1},\ldots,1)$ makes $U_zv=bv$ for a nonzero
$b$ (the empty product is $1$ when $n=1$).  Thus
$v=U_z^m(b^{-m}v)$ for all $m$.  AFS11 implies that its
Jacquet image $x=p(v)$ is nonzero.

Let $J$ be the normalized Borel Jacquet module.  It has finite
length by the retained Bernstein--Zelevinsky theorem.  As
explained in AFS12, every smooth irreducible representation
of the diagonal torus is a character: an open compact
stabilizer of a cyclic vector acts trivially, and the
remaining finitely generated abelian quotient has
one-dimensional simple complex modules by the weak
Nullstellensatz.  Hence $J$ is finite-dimensional.
The finite-dimensional commutative algebra generated by the
torus action decomposes $J$ into generalized-character blocks.
Choose one block with a nonzero projection of $x$, and let
its character be $\mu$.  If $t x=cx$, projection to this
block gives the same scalar relation for a nonzero vector;
there $t-\mu(t)$ is nilpotent, so necessarily $c=\mu(t)$.
The chosen block has a nonzero character quotient: quotient
by its maximal ideal and use Nakayama, then a nonzero linear
functional.  Extend by zero on the other blocks.  This gives
a nonzero functional $l:J\to\mathbb C$ with
$l(tu)=\mu(t)l(u)$.  It need not be nonzero on $x$; the
scalar-relation property just proved is what is required.
The character is smooth because a vector on which $l$ is
nonzero has an open stabilizer.

Write $\mu(t)=\prod_i\eta_i(t_i)$.  On compact torus units
the Jacquet modulus is one, so the scalar relations of $x$
give $\eta_i(u)=\varphi(e_i(u))^{-1}$.  The explicit
Frobenius map
\[
        v'\longmapsto [g\longmapsto l(p(\sigma(g)v'))]
 \tag{HAW.12}
\]
is nonzero, smooth, equivariant, and has left $B$-law
$\delta_B^{1/2}\mu$.  Compactness of $B\backslash G$
ensures it lies in the original induced representation.
Irreducibility makes it an injection.  Put
$\omega_i=\eta_i|\cdot|^{-(n-1)/2}$; the target is exactly
$\mathrm{PS}(\omega)$ and the unit relation is unchanged.

It remains to keep the scalar normalization.  The actual
upper-unipotent coset index is $I_j=q^{j(n-j)}$.  AFS11's
unnormalized Jacquet identity and normalized torus action give
\[
 p(U_{a_j}v)=I_jp(\sigma(a_j)v),\qquad
 \mu(a_j)=\delta_B(a_j)^{-1/2}I_j^{-1}b_j.
\]
Since $\delta_B(a_j)=q^{-j(n-j)}$, we obtain
\[
 b_j=I_j\delta_B(a_j)^{1/2}\mu(a_j)
     =q^{j(n-j)/2-j(n-1)/2}\prod_{i<j}\omega_i(\varpi),
\]
which is (HAW.11), also for $j=n$.  This proves the full
original $\mathrm{OrdJac}$ without distinct unit characters
or a good-point hypothesis.
\end{proof}
% END HAW7

% BEGIN HAW8
\paragraph{HAW8. The original scalar extraction OrdScalAlt.}
\label{claim:HAW8}
Retain all data and guards in $\mathrm{OrdScalAlt}$:
$H.\mathrm{Ax}$, $\mathrm{IsCO}(O,A)$, $H.\mathrm{IsWt}(\theta)$,
an actual $\mathrm{IsSpec}(\theta,\iota,x,\lambda,\xi,c_0,sp)$,
an original part $Y$, a level $K$, and $e\ge c_0$.
If $0\ne F\in Y.Y\cap\mathrm{ClAt}(K,e)$ is an eigenvector
of every local positive generator at $v$, with nonzero
eigenvalues $b_j$, there are smooth $\omega_i$ with
\[
 \begin{split}
 Y.\sigma_l(v)&\text{ a subquotient of }\mathrm{PS}(\omega),\\
 \omega_i(u)\,\iota(\mathrm{finPart}(\theta,x,\lambda,v,i,u))&=1,\\
 b_j&=\mathrm{lowVal}(\iota,\lambda,v,j+1)^{-1}
                         \mathrm{psEig}(\omega,j+1).
 \end{split}
 \tag{HAW.13}
\]
\begin{proof}
The original $\mathrm{ClAt}$ conditions make $F$ fixed by
the local upper unipotent and the lower depth-$e$ subgroup.
The compact torus acts through the inverse of its finite-part
values.  These values give an actual character: uniqueness of
a scalar acting on the nonzero $F$ makes the identity and
multiplication laws follow from the representation law, and
invertibility makes every scalar nonzero.  The specialization's
triviality condition makes the character trivial on
$U_c(k,e+1)$, since $e\ge c_0$.

The local isotypy $Y.\mathrm{isoL}(v)$ and AFS13 supply an
equivariant linear map to $Y.\sigma_l(v)$ nonzero on $F$.
All invariance and scalar relations survive this map.  There
is no tensor-product completion or scalar extension of the
chosen part in this operation.  The global operator supported
at $v$ is the actual local operator: its coset representatives
are obtained by putting each local representative in the $v$
coordinate and $1$ in every other coordinate.  This gives
a bijection of the original cosets, without extra index.
Thus on $F$,
\[
 \mathrm{opC}_{v,j}F
       =\mathrm{lowVal}_{v,j+1}^{-1}U_{a_{j+1}}F.
\]
Every displayed low value is nonzero.  If it were zero, the
source convention for inversion would make the left side
zero, contradicting $b_jF\ne0$.  Set
$a_{j+1}=\mathrm{lowVal}_{v,j+1}b_j\ne0$ and apply HAW7 to
the nonzero local image of $F$.  Dividing its resulting
$\mathrm{psEig}$ formula by that same low value gives
(HAW.13).  These are precisely all fields of
$\mathrm{OrdScalAlt}$, for the supplied data and every $v$.
\end{proof}
% END HAW8

% BEGIN HAW9
\paragraph{HAW9. PartPtAlt for an arbitrary predicate on the actual part.}
\label{claim:HAW9}
Suppose the original $\mathrm{AutYAlt}(H,\iota,P)$ holds,
where $P(\lambda,d,\xi,Y)$ may depend on the entire part.
At every specified arithmetic specialization and every
specified $O$-algebra point $y$ of the original commuting
Hecke algebra extending $x$, there are an actual part $Y$
and $F\ne0$ such that
\[
 P(\lambda,d,\xi,Y),\qquad
 F\in Y.Y\cap\mathrm{OrdAt}(K,c_0),\qquad
 \mathrm{opC}_iF=\iota(y(\mathrm{gen}_i))F\quad\text{for all }i.
 \tag{HAW.14}
\]
The hypotheses are exactly those of $\mathrm{PartPtAlt}$,
including the specified $\mathrm{IsSpec}$ map and admitted
weight; no property of $P$ beyond $\mathrm{AutYAlt}$ is used.
\begin{proof}
CFD9--CFD12's finite averaging and polynomial generalized-eigenprojection
argument applies verbatim to the collection of parts satisfying
$P$.  Indeed $\mathrm{AutYAlt}$ supplies the needed spanning
statement for every actual classical vector; every part has
the same stability fields, and AFS16 preserves each under the
original Hecke operators.  Average over the original level,
upper and lower compact subgroups and the compact torus
character to preserve each part.  These are finite averages
over cosets of stabilizers; normality of those stabilizers is
not required, and translating cosets proves invariance.
On the resulting finite-dimensional classical space, use the
original normalized distinguished $z$-operator.  By CFD11,
the ordinary space is exactly the sum of its generalized
eigenspaces whose eigenvalues $\alpha$ satisfy
$|\iota^{-1}(\alpha)|_\ell=1$.  The Chinese remainder
theorem supplies a polynomial equal to $1$ modulo the
primary factors for precisely those eigenvalues and $0$
modulo every other primary factor, including any nonzero
nonunit eigenvalue.  This polynomial in the original
operator preserves each part and projects exactly onto
the original ordinary space.  Therefore
\[
       \mathrm{OrdAt}=\sum_{Y:P(Y)}(Y.Y\cap\mathrm{OrdAt}).
\]
ASP11 identifies this space with the complex scalar extension
of the actual finite free ordinary module; ASP13 intertwines
all the original Hecke generators.  The residual idempotent
of OHA6--OHA7 selects its specified ordinary block and
preserves every part by AFS18.  The Hecke algebra acts faithfully on
that block by its definition as a subalgebra of endomorphisms.

AFS9 now applies to the given point $y$: its module fiber is
nonzero, and the corresponding local Artinian block of the
finite commuting algebra is nonzero.  The sum above contains
a nonzero contribution in that block from some $Y$ with $P$.
In a last nonzero power of the block's nilpotent maximal
ideal, a nonzero vector is a simultaneous eigenvector for
the given character.  The specialization map transports it
injectively to the required $F$, and intertwining gives every
equation in (HAW.14).  This proves $\mathrm{PartPtAlt}$ for
the full part-dependent predicate, not only predicates of
the frame representation.
\end{proof}
% END HAW9

% BEGIN HAW10
\paragraph{HAW10. Full CoreOrdAlt with no good-point restriction.}
\label{claim:HAW10}
Retain exactly the $\mathrm{CoreOrdAlt}$ data in
\texttt{OrdDefs}: an original $H,A,\theta$ with its axioms,
an arbitrary $\mathrm{AutYAlt}$ predicate $P$, the specified
level and residual Hecke algebra $T$, its $\mathrm{IsCO}$
and commutativity structures, and the actual
$\mathrm{IsUniv}(D,\psi)$ record.  For every admitted
$\lambda$ and every $y:T\to_O\overline{\mathbb Q}_\ell$
arithmetic of weight $\lambda$, there exist $d,\xi,Y$ with
\[
 \begin{split}
 &P(\lambda,d,\xi,Y),\qquad D\otimes y=\det Y.r,\\
 &\forall v\ \exists\omega:\quad
 Y.\sigma_l(v)\text{ is a subquotient of }\mathrm{PS}(\omega),\\
 &\hspace{30mm}\forall i\quad
 \mathrm{OrdChar}(\iota,H.\mathrm{art}(v),\omega_i,
                       \mathrm{htW}(\lambda,i),y\psi_{v,i}).
 \end{split}
 \tag{HAW.15}
\]
The determinant equality is the full polynomial law and the
ordinary characters have the exact original order.
\begin{proof}
Put $x=y|_A$.  ASP13 gives the original arithmetic
specialization, and HAW9 gives $Y,F$ with all the required
eigenvalue equations.  HEP7's determinant argument gives
$D\otimes y=\det Y.r$; it compares the original spherical
eigenvalues and then the full characteristic-zero determinant
law, including arbitrary scalar algebras.

For each local generator, $\mathrm{IsUniv}.\psi\_unif$
expresses its value as a product of values of the unit-valued
characters $\psi_{v,i}$ at an Artin preimage of a uniformizer.
Thus $y(\mathrm{gen}_{v,j})\ne0$, and so does its image
under $\iota$.  Membership in $\mathrm{OrdAt}$ implies
membership in the corresponding $\mathrm{ClAt}$, so HAW8
applies to these nonzero eigenvalues.  Its inverse-unit
relation and exact $\mathrm{psEig}$ equality are precisely
the hypotheses of HEP7's local ordinary-character statement.
That statement gives the final line of (HAW.15), including
continuity of each $y\psi_{v,i}$.  All fields refer to the
same $Y,D,\psi,y$.  No pairwise distinct finite characters,
goodness condition, or flag on $Y.r$ was used.
\end{proof}
% END HAW10

% BEGIN HAW11
\paragraph{HAW11. EndOrdAlt for the same arbitrary representing rho.}
\label{claim:HAW11}
For every original $R.\mathrm{IsJ}(B,K,J)$, the full
$\mathrm{ENDO}(B,\iota,J)$ holds.  Thus if $J\subseteq\ker x$,
$x.\mathrm{Arith}(h)$, $x.\mathrm{IsRep}(\rho)$ and $\rho$
is irreducible, there is an actual regular algebraic cusp
$\pi$ with $\mathrm{DualIsConj}(\pi)$ and
$\mathrm{AttC}(\iota,\rho,\pi)$ such that, at every chosen
$w=s.w(v)$ over $\ell$, $\mathrm{ExpOf}(\iota,s.hw(v),
\pi,x.\psi_v)$ holds.  Its exact witnesses are an original
local Artin map, smooth $\omega_i$, and integers $h_i(\tau)$
strictly increasing in $i$ for every embedding, with
\[
 \mathrm{OrdChar}(\iota,a,\omega_i,h_i,x.\psi_{v,i}),\qquad
 \forall\tau_w\text{ a local component of }\pi,\quad
                 \tau_w\text{ a subquotient of }\mathrm{PS}(\omega).
 \tag{HAW.16}
\]
Continuity of the supplied $\rho$ is not a premise.
\begin{proof}
The $J=\top$ branch admits no such $x$, because $1$ would
map to zero in a field.  In the nonzero branch retain the
original surjection $f:B\twoheadrightarrow T$, its
$\mathrm{IsMap}$ record, and its compatibility with the
given weight algebra.  The map $x$ factors uniquely as
$y\circ f$, since $\ker f=J\subseteq\ker x$.
AFS6--AFS8 convert the original arithmetic weight to the
admitted dominant $\lambda$ for this same $y$.
Apply HAW10 using HAW2's $\mathrm{AutUOAlt}$ predicate.
Naturality of scalar extension and pushing characters gives
$x.D=\det Y.r$ and exactly the characters $x.\psi$ in
(HAW.15).  Pullback along $V.\mathrm{pr}$ gives
$x.\mathrm{IsRep}(Y.r\circ V.\mathrm{pr})$ as a full law.

The map $V.\mathrm{pr}$ is surjective: lift a frame element
to the original global quotient, then use $V.h\eta$ to lift
it to $G_L$.  An invariant subspace for the pullback is
therefore precisely an invariant subspace for $Y.r$.
Consequently the pullback is semisimple.  The given
irreducible $\rho$ is semisimple too.  Equality of their
determinant laws gives equality of their characteristic
polynomials, and the finite algebra argument in HAW1 gives
an isomorphism, without requiring continuity of $\rho$.
Thus the pullback is irreducible.  The part predicate supplies
$\pi$, and attachment transfers through this isomorphism.

For (HAW.16) keep the Artin map and $\omega$ in (HAW.15)
and take
$h_i(\tau)=\mathrm{htW}(\lambda,i)(\tau)
=\lambda_\tau(n-1-i)+i$.
Dominance makes the first term nondecreasing in $i$, while
$i$ is strictly increasing; hence these integer weights are
strictly increasing.  HAW2 identifies $Y.\sigma_l(v)$ with
an actual component of $\pi$.  Every other component at
that place is isomorphic to it, so the same
principal-series subquotient statement holds for each one.
This proves the full $\mathrm{ExpOf}$ and hence the original
$\mathrm{EndOrdAlt}$, for the given $\rho$ and same ideal.
\end{proof}
% END HAW11

\subsection{Ordered flags and transport between conjugate places}

% BEGIN HAW12
\paragraph{HAW12. Ordered matrix relations force the original flag.}
\label{claim:HAW12}
Let $K$ be a field, $\rho:H\to\mathrm{GL}_n(K)$, and
$\psi_i:H\to K^\times$.  Suppose
\[
 \prod_{i=0}^{n-1}(\rho(g_i)-\psi_i(g_i)I)=0,\qquad
 \det(X-\rho(g))=\prod_i(X-\psi_i(g)),
 \tag{HAW.17}
\]
for all ordered tuples and all $g$, and some $s\in H$
has pairwise distinct $\psi_i(s)$.  Then $\rho$ has the
full upper-triangular flag with ordered diagonal characters
$\psi_0,\ldots,\psi_{n-1}$.  For the original arithmetic
characters $x.\psi_v$, such a separating $s$ always exists.
\begin{proof}
Put $\lambda_i=\psi_i(s)$.  The characteristic polynomial
of $\rho(s)$ has distinct roots, so its Lagrange spectral
idempotents $E_i$ are nonzero rank-one projections whose sum
is $I$.  In (HAW.17) set every $g_j=s$ except $g_b=g$,
and multiply on the left by $E_a$ and on the right by $E_b$.
For $a\ge b$ this gives
\[
 \Bigl(\prod_{j<b}(\lambda_a-\lambda_j)\Bigr)
 E_a(\rho(g)-\psi_b(g)I)E_b
 \Bigl(\prod_{j>b}(\lambda_b-\lambda_j)\Bigr)=0.
 \tag{HAW.18}
\]
Both scalar factors are nonzero.  For $a>b$ this says
$E_a\rho(g)E_b=0$; for $a=b$ it says
$E_b\rho(g)E_b=\psi_b(g)E_b$.  Choosing one nonzero vector
in each $E_iK^n$, in this order, proves the exact flag.
No factors from different positions have been commuted.

At a place over $\ell$, arithmeticity expresses each
$\psi_i$ on inertia as a smooth finite-part character times
the algebraic character $\prod_\tau\tau(u)^{-h_i(\tau)}$.
There are finitely many local embeddings and at least one,
since the completion is a finite extension of $\mathbb Q_\ell$.
Intersect the finitely many open kernels of the finite parts
in the compact inertia group; its normal subgroup $U$ has
finite positive index $e$.  Local reciprocity lifts the
rational unit $1+\ell$ to an inertia element $\gamma$.
Then $\gamma^e\in U$, and
\[
       \psi_i(\gamma^e)
                  =((1+\ell)^e)^{-\sum_\tau h_i(\tau)}.
 \tag{HAW.19}
\]
The weights strictly increase at each embedding, so their
sums strictly increase.  Distinct integer powers of the
positive rational number $(1+\ell)^e>1$ remain distinct
in $\overline{\mathbb Q}_\ell$.  Thus $s=\gamma^e$
separates the original characters.  This uses arithmeticity,
not a good-point condition.
\end{proof}
% END HAW12

% BEGIN HAW13
\paragraph{HAW13. The full original OrdFlagJ for arbitrary irreducible rho.}
\label{claim:HAW13}
Let $B.\mathrm{RingProps}(\Lambda,J)$ hold.  If
$J\subseteq\ker x$, $x.\mathrm{Arith}(h)$,
$x.\mathrm{IsRep}(\rho)$ and $\rho$ is irreducible, then
at each chosen place over $\ell$ its restriction has
$\mathrm{HasFlag}$ with the \emph{ordered} characters
$x.\psi_{v,i}$.  This is the full $\mathrm{FlagJ}$ predicate;
there is no continuity hypothesis on $\rho$ and no
$\mathrm{V0ok}$ guard.
\begin{proof}
Irreducibility implies $n\ge1$; the zero-dimensional vector
space cannot be an irreducible nonzero module.  The original
$\phi=V.\mathrm{pr}$ is surjective: lift a frame element to
the original global quotient and apply $V.h\eta$ to obtain
its preimage in $G_L$.
We first descend $\rho$ to the frame.  If $g\in\ker\phi$,
the determinant identity $x.D\circ\phi=\det\rho$ gives
\[
          \operatorname{tr}(\rho(t)(\rho(g)-I))=0
                                           \quad(t\in G_L).
\]
The matrix span of an irreducible representation over an
algebraically closed field is all $M_n$ (the finite algebra
density theorem).  The trace pairing on $M_n$ is
nondegenerate, as is seen by testing matrix units.  Hence
$\rho(g)=I$.  Thus $\rho=\rho_0\circ\phi$, and $\rho_0$
has the same full matrix span.  Its determinant equals
$x.D$ as a full polynomial law: for any scalar algebra,
surjectivity of $\phi$ makes the induced map of group
algebras surjective, so equality after pullback implies
equality before it.

Factor $x$ through $B/J$.  Specializing $\mathrm{RingProps}.O1$
gives the characteristic polynomial identity in (HAW.17)
for the local restriction of $\rho_0$.  Specializing its
$O2$ gives the literal $\mathrm{Kills}$ assertion for the
ordered group-algebra element
$z=\prod_i([g_i]-x.\psi_i(g_i)[1])$.
To spell out its implication, use the scalar algebra
$\overline{\mathbb Q}_\ell[T]$ and test $\mathrm{Kills}$
on $T[t]$.  The coefficient of $T$ in
$\det(I+T\rho_0(t)\rho_0(z))=1$ is
$\operatorname{tr}(\rho_0(t)\rho_0(z))=0$.
The full matrix span and nondegenerate trace pairing force
$\rho_0(z)=0$.  This is the ordered product identity in
(HAW.17), with the exact source order.

HAW12 provides the separating local element and converts
these two identities into the required flag.  Pullback to
the original local decomposition group is exactly the
restriction of $\rho$ in the assertion.  Only $O1,O2$ of
the ring record were used; the conditional auxiliary
unramified field plays no role.
\end{proof}
% END HAW13

% BEGIN HAW14
\paragraph{HAW14. Polarization and full ConjUnl transport.}
\label{claim:HAW14}
Let $\rho$ be continuous, let $\pi$ satisfy
$\mathrm{DualIsConj}(\pi)$ and $\mathrm{AttC}(\iota,\rho,\pi)$,
and let $w,w'$ lie over the same real place.  Fix one
prolongation $A$ over $w$ and a full flag of $\rho|_{D_A}$
with characters $\psi_i$.  Then
\[
 \begin{split}
 \mathrm{Unl1}(\rho,w')&\Longrightarrow
                         \forall i,j\quad\psi_i\ne\psi_j\varepsilon,\\
 [\forall i,j, m\ge1\quad\psi_i^m\ne(\psi_j\varepsilon)^m]
                 &\Longrightarrow\mathrm{PotUnl}(\rho,w').
 \end{split}
 \tag{HAW.20}
\]
These are both clauses of the literal $\mathrm{ConjUnl}$.
\begin{proof}
We first record the polarization consequence used in the
transport.  For every complex-conjugation lift $c$, the
continuous semisimple representations
$\rho^c$ and $\rho^\vee\varepsilon^{1-n}$ are isomorphic.
Here semisimplicity is included in $\mathrm{AttC}$.
Indeed the original duality dictionary (finite automorphic
type and Gelfand--Kazhdan) converts $\mathrm{DualIsConj}$
to the corresponding conjugate-self-dual cuspidal class.
At cofinally many places, it sends the normalized Satake
multiset to its inverse at the conjugate place.  The attached
geometric Frobenius roots are $q_w^{(n-1)/2}\alpha$.
Dualizing and twisting by $\varepsilon^{1-n}$ sends them to
\[
       (q_w^{(n-1)/2}\alpha)^{-1}q_w^{n-1}
                             =q_w^{(n-1)/2}\alpha^{-1}.
\]
Intersect the attachment set with its conjugate to obtain
the same characteristic polynomials for the two
representations.  Chebotarev and the finite semisimple
comparison in HAW1 give their isomorphism.  Equivalently it
is a nondegenerate pairing with the original multiplier
$\varepsilon^{1-n}$.  No symmetry sign of that pairing is
needed for this consumer.  This reconstructs the finite
polarization comparison from the retained original
automorphic conjugation and Frobenius laws.

For any prolongation $A'$ over $w'$, transitivity of the
absolute Galois action on prolongations over the common real
place gives $g\in G_{L^+}$ carrying $A$ to $A'$.  Conjugation
gives an isomorphism $f:D_{A'}\simeq D_A$ and preserves
the cyclotomic character.  The restriction of $g$ to $L$ is
either the identity or complex conjugation.  In the first
case, the two restrictions of $\rho$ are conjugate after
precomposition by $f$.  In the second, write $g$ as an
element of $G_L$ times a complex-conjugation lift and use
the preceding pairing.  The second restriction is then
conjugate to
\[
          (\rho|_{D_A})^\vee\circ f\ \otimes\mu
 \tag{HAW.21}
\]
for the common character $\mu=\varepsilon^{1-n}$ pulled
through $f$.  These statements use actual decomposition
groups, not just Weil semisimplifications.

A flag transports under a group isomorphism and a conjugacy.
The dual flag has the reverse-ordered inverse characters;
a common twist multiplies all of them by $\mu$.
For each $m\ge1$, existence of a cyclotomic ratio is unchanged.
In the dual case the exact equivalence, with the pair reversed,
is
\[
 \psi_i^m=(\psi_j\varepsilon)^m
 \quad\Longleftrightarrow\quad
 (\psi_j^{-1}\mu)^m=(\psi_i^{-1}\mu\varepsilon)^m.
 \tag{HAW.22}
\]
Surjectivity of $f$ gives the same equivalence before and
after precomposition.

Finally two full flags of the same representation have the
same multiset of characters.  Their traces are the sums of
their diagonal characters.  Distinct group characters are
linearly independent: in a shortest nonzero relation,
translate by an element on which two characters differ and
subtract a scalar multiple of the original relation to
obtain a shorter one.  Thus their multiplicities agree in
characteristic zero.  This argument does not require local
semisimplicity.  If $\mathrm{Unl1}(\rho,w')$ holds, compare
its flag at any $A'$ with the transported flag and use
(HAW.22) for $m=1$; a forbidden pair among the $\psi_i$
would be a forbidden pair in that flag.  If all powers are
excluded among the $\psi_i$, (HAW.22) excludes them in the
transported flag at \emph{every} $A'$.  This is exactly
$\mathrm{PotUnl}$ and proves (HAW.20).
\end{proof}
% END HAW14

\subsection{The exact strict window excludes every torsion ratio}

% BEGIN HAW15
\paragraph{HAW15. The ordinary-character linking calculation.}
\label{claim:HAW15}
Let $w\mid\ell$, fix the original local Artin map
$a:W_{L_w}\to L_w^\times$, and suppose
$\mathrm{OrdChar}(\iota,a,\omega_i,h_i,\psi_i)$ holds,
with $h_i(\tau)$ strictly increasing in $i$ at every
embedding.  Then
\[
          \omega_i=\omega_j|\cdot|_w
                      \quad\Longrightarrow\quad
                         \psi_i=\psi_j\varepsilon.
 \tag{HAW.23}
\]
More precisely, put
$D(\tau)=h_j(\tau)-h_i(\tau)-1$ and
$\delta=\psi_i/(\psi_j\varepsilon)$.  For every Weil
element $s$, the exact identity is
\[
 \delta(s)\,\iota^{-1}(\omega_j(a(s)))q_w^{\deg s}
       =\iota^{-1}(\omega_i(a(s)))
                           \prod_\tau\tau(a(s))^{D(\tau)}.
 \tag{HAW.24}
\]
\begin{proof}
The defining ordinary-character identity is
$\psi_k(s)=\iota^{-1}(\omega_k(a(s)))
\prod_\tau\tau(a(s))^{-h_k(\tau)}$.
Local reciprocity in the original geometric Artin convention
gives
\[
 \varepsilon(s)=q_w^{\deg s}\prod_\tau\tau(a(s)),\qquad
 |a(s)|_w=q_w^{\deg s}.
 \tag{HAW.25}
\]
Divide the identities, keeping the same local embeddings;
all factors are nonzero.  This proves (HAW.24).

A continuous character of a compact group into
$\overline{\mathbb Q}_\ell^\times$ has $\ell$-adic norm
one: otherwise positive powers of it or of its inverse
would be unbounded, contrary to compactness of its image.
The $\psi_i$ are continuous by $\mathrm{OrdChar}$, hence
so is $\delta$ on the compact decomposition group.
Under the premise in (HAW.23), equations (HAW.24)--(HAW.25)
reduce on the Weil group to
$\delta(s)=\prod_\tau\tau(a(s))^{D(\tau)}$.
Choose an Artin preimage of the rational nonzero number
$\ell$.  Its value is $\ell^{\sum_\tau D(\tau)}$,
whose $\ell$-adic norm can be one only if
$\sum_\tau D(\tau)=0$.

The strictly increasing integer weights now force
$D(\tau)=0$ for every local embedding.  If $i<j$, every
$D(\tau)\ge0$ and their sum is zero.  If $i\ge j$,
every $D(\tau)<0$, which is impossible since the set of
local embeddings is nonempty.  Thus $\delta=1$ on the
Weil group.  This subgroup is dense in the decomposition
group, and $\delta$ is continuous with Hausdorff target,
so $\delta=1$ on the entire decomposition group.  This
proves (HAW.23) with no extra continuity assumption on
any ambient global representation.
\end{proof}
% END HAW15

% BEGIN HAW16
\paragraph{HAW16. Torsion ratios have exactly the forbidden factor q; full ExpTor.}
\label{claim:HAW16}
Under HAW15's hypotheses, for every $m\ge1$,
\[
 \psi_i^m=(\psi_j\varepsilon)^m
       \quad\Longrightarrow\quad
                   |\omega_j(\varpi_w)|=q_w|\omega_i(\varpi_w)|.
 \tag{HAW.26}
\]
Consequently the full original $\mathrm{ExpTor}$ holds:
if $\mathrm{ExpOf}(\iota,a,\pi,\psi)$ holds at $w\mid\ell$,
then $\mathrm{TorAt}(\pi,w)$; if also
$\psi_i\ne\psi_j\varepsilon$ for all $i,j$, then
$\mathrm{TU}(\pi,w)$; and any $\mathrm{WinAt}(\pi,t,w)$
implies the exclusions (HAW.3) for these same $\psi_i$.
\begin{proof}
Suppose $\delta^m=1$ with $\delta$ as in HAW15.
Smoothness of $\omega_i,\omega_j$ makes the intersection
of their kernels an open neighborhood of $1$ in
$L_w^\times$.  The rational units $1+\ell^M$ converge to
$1$, so choose $M\ge1$ with $N=1+\ell^M$ in that
intersection.  Here $N>1$ and its local valuation is zero.
An Artin preimage $s$ of $N$ has degree zero.
Equation (HAW.24) gives $\delta(s)=N^{\sum_\tau D(\tau)}$.
Raising to $m$ shows $N^{m\sum D}=1$ in
$\overline{\mathbb Q}_\ell$, hence in $\mathbb Q$.
Since $N>1$ and $m>0$, $\sum D=0$.  The integer-weight
argument in HAW15 again gives $D(\tau)=0$ at every local
embedding.

Take now an Artin preimage of $\varpi_w$.  Its degree is
$-1$, so (HAW.24) becomes
$\delta(s)\iota^{-1}(\omega_j(\varpi_w))q_w^{-1}
=\iota^{-1}(\omega_i(\varpi_w))$.
The complex image of $\delta(s)$ is an $m$-th root of unity
and has modulus one.  Taking complex absolute values proves
(HAW.26), with the exact factor $q_w$.

For the three $\mathrm{ExpTor}$ conclusions, take its
specified smooth characters $\omega$ and any actual local
component of $\pi$.  It is a subquotient of
$\mathrm{PS}(\omega)$, whose normalized inducing characters
are
\[
             \mu_i=\omega_i|\cdot|_w^{(n-1)/2}.
 \tag{HAW.27}
\]
Cuspidal support uniqueness identifies the support of every
component, and every support witness, with the multiset of
these rank-one characters.  This proves $\mathrm{TorAt}$.
If (HAW.1) failed for $\mu_i,\mu_j$, their common shift
cancels, and the unramified quotient of value $q_w$ gives
$\omega_j=\omega_i|\cdot|_w$.  Formula (HAW.23), with
indices interchanged, would give
$\psi_j=\psi_i\varepsilon$.  Thus the stated exclusion
implies $\mathrm{TU}$.

Finally suppose $\mathrm{WinAt}(\pi,t,w)$ and a torsion
ratio existed.  Formula (HAW.26) remains true for $\mu_i,
\mu_j$, since the shift in (HAW.27) is common.  The two
strict window inequalities would give
\[
 |\mu_j(\varpi_w)|=q_w|\mu_i(\varpi_w)|
       >q_wq_w^{t-1/2}=q_w^{t+1/2},
\]
contradicting the strict upper bound for $\mu_j$.
This proves the exclusions simultaneously for every pair
and every positive integer $m$.  There is no bound on $m$
and no loss of either endpoint of the window.
\end{proof}
% END HAW16

% BEGIN HAW17
\paragraph{HAW17. The full BornJ implication above ell.}
\label{claim:HAW17}
Assume the original $\mathrm{RingProps}(\Lambda,J)$ and
$\mathrm{ENDO}(B,\iota,J)$.  Let $x,h,\rho$ satisfy
$J\subseteq\ker x$, $x.\mathrm{Arith}(h)$,
$x.\mathrm{IsRep}(\rho)$, irreducibility, and continuity
of $\rho$.  If $w\mid\ell$ and
$\mathrm{Unl1}(\rho,w)$, then
\[
 \mathrm{PotUnl}(\rho,w),\qquad
 \forall\pi\quad\mathrm{AttC}(\iota,\rho,\pi)
                              \Longrightarrow\mathrm{TU}(\pi,w).
 \tag{HAW.28}
\]
This is the entire original $\mathrm{BornJ}$, with all
prolongations and every attached cusp in its second clause.
\begin{proof}
Use $\mathrm{ENDO}$ to obtain one actual cusp $\pi_0$,
its dual-is-conjugate property, attachment, and all its
$\mathrm{ExpOf}$ witnesses.  Put
$v=\underline w$ and $w_0=s.w(v)$.  The original place
dictionary says $w_0\mid\ell$ and
$\underline{w_0}=\underline w$.  HAW13 gives the ordered
flag $\psi=x.\psi_v$ at the chosen prolongation over $w_0$.
The first clause of HAW14 turns
$\mathrm{Unl1}(\rho,w)$ into
$\psi_i\ne\psi_j\varepsilon$.  HAW16 then makes
$\mathrm{TU}(\pi_0,w_0)$.

Choose the single global window center for $\pi_0$ by
HAW6.  Its window at $w_0$ and the last clause of HAW16
exclude every positive-power cyclotomic ratio among the
same $\psi_i$.  The second clause of HAW14 therefore gives
$\mathrm{PotUnl}(\rho,w)$ at every prolongation.

Let $\pi$ be any other attached cusp.  At cofinally many
places its Satake multiset and that of $\pi_0$ give the
same Frobenius polynomial for $\rho$.  Equality of their
root multisets, followed by cancellation of the common
nonzero scalar $q^{(n-1)/2}$, gives the same normalized
Satake multisets with multiplicities.  The original strong
multiplicity-one theorem identifies their cuspidal classes;
its actual-space version identifies the corresponding
$\mathrm{CuspSub}$, and in either formulation every local
component and support is the same up to isomorphism.
HAW4 transports $\mathrm{TU}(\pi_0,w_0)$ to
$\mathrm{TU}(\pi_0,w)$, and this identification gives
$\mathrm{TU}(\pi,w)$.  No regular-algebraic guard was
needed for this second attached cusp.
\end{proof}
% END HAW17

% BEGIN HAW18
\paragraph{HAW18. The same frame point with the full AwayOk predicate.}
\label{claim:HAW18}
If $V.S'=s.S$ and $R.\mathrm{IsJ}(B,K,J)$, every
$x,h$ with $J\subseteq\ker x$ and $x.\mathrm{Arith}(h)$
has a continuous semisimple frame representation
$\rho_0:V.\mathrm{frame}.G_0\to\mathrm{GL}_n
(\overline{\mathbb Q}_\ell)$ with
\[
 x.\mathrm{IsRep}(\rho_0\circ V.\mathrm{pr}),\qquad
 \mathrm{IsIrr}(\rho_0\circ V.\mathrm{pr})
       \Longrightarrow\mathrm{AwayOk}(s,\iota,
                                     \rho_0\circ V.\mathrm{pr}).
 \tag{HAW.29}
\]
This is the literal $\mathrm{coreA}$ output, including its
predicate and original frame.
\begin{proof}
Again $J=\top$ is impossible at the given point.  Retain
the actual surjection $B\twoheadrightarrow T$, factor
$x$ through it, and use AFS6--AFS8 to obtain the same
admitted dominant weight.  HAW5 is an
$\mathrm{AutomorphicHeckeSet}$ spanning assertion.  Viewed
as a predicate of the full part via its frame representation,
it is exactly the $\mathrm{AutYAlt}$ hypothesis in HAW10.
Apply HAW10 to that predicate and the given point of $T$.
Put $\rho_0=Y.r$ for the resulting part.  Continuity and
semisimplicity are fields of that original part.  Its
predicate gives the implication in (HAW.29), while its full
determinant law, composed with the original surjection and
$V.\mathrm{pr}$, gives the first assertion.  Thus this
construction also proves the stronger point predicate
without changing $T$, the point, the frame, or the ideal.
\end{proof}
% END HAW18

% BEGIN HAW19
\paragraph{HAW19. The full literal IrrArithPointsAutomorphicAway.}
\label{claim:HAW19}
Under exactly the original $H'$, for every original datum
with $n\ge1$, $V.S'=s.S$ and $R.\mathrm{IsJ}(B,K,J)$,
\[
              \mathrm{IrrArithPointsAutomorphicAway}(B,\iota,J)
 \tag{HAW.30}
\]
holds.  Explicitly, for every $x,h,\rho$ with
$J\subseteq\ker x$, $x.\mathrm{Arith}(h)$,
$x.\mathrm{IsRep}(\rho)$ and $\rho$ irreducible, the full
$\mathrm{AutOffN}(s.v_0,\iota,\rho)$ holds.  Its clauses are
\[
 \begin{split}
 &\mathrm{AutCsd0}(\iota,\rho),\\
 &\mathrm{Continuous}(\rho)\ \Longrightarrow\\
 &\quad\left\{\begin{array}{l}
     \forall w\mid\ell,\quad
       \mathrm{Unl1}(\rho,w)\Longrightarrow\mathrm{PotUnl}(\rho,w),\\
     \forall\pi,\ \mathrm{RegAlg}(\pi)\land
                  \mathrm{AttC}(\iota,\rho,\pi)\Longrightarrow\\
     \qquad\exists t\in\mathbb R\ \forall w,\quad
     \underline w\ne s.v_0\ \land\
       (w\nmid\ell\ \lor\ \mathrm{PotUnl}(\rho,w))
          \Longrightarrow\mathrm{WinAt}(\pi,t,w).
   \end{array}\right.
 \end{split}
 \tag{HAW.31}
\]
\begin{proof}
HAW18 provides $\rho_0$ with the same determinant law.
Surjectivity of $V.\mathrm{pr}$ makes its pullback
semisimple, as in HAW11.  The given $\rho$ is irreducible
and hence semisimple.  HAW1's finite algebra comparison
therefore gives
$\rho_0\circ V.\mathrm{pr}\simeq\rho$, without a
continuity premise on the latter.  Transport irreducibility
to the former and apply its $\mathrm{AwayOk}$ predicate.
Its regular algebraic dual-is-conjugate cusp $\pi_0$ is
attached also to $\rho$ by isomorphism.  This is exactly
$\mathrm{AutCsd0}$.

To prove the second clause, now assume
$\mathrm{Continuous}(\rho)$, exactly where the source
places this implication.  JRP15 supplies the full
$\mathrm{RingProps}(\Lambda,J)$ at this same ideal, and
HAW11 supplies $\mathrm{ENDO}$.  Apply HAW17.  It gives
the required $\mathrm{Unl1}\Rightarrow\mathrm{PotUnl}$
at every place above $\ell$.

Fix any regular algebraic attached cusp $\pi$.  Choose
$t$ once, before choosing a place, from HAW6 applied to
this $\pi$.  Let $w$ satisfy the two guards in (HAW.31).
If $w\nmid\ell$, the $\mathrm{AwayOk}$ property gives
$\mathrm{TU}(\pi_0,w)$.  The same strong-multiplicity-one
comparison as in HAW17 transports it to
$\mathrm{TU}(\pi,w)$.  If $w\mid\ell$, the other guard
is $\mathrm{PotUnl}(\rho,w)$.  Its exclusions for $m=1$
give $\mathrm{Unl1}(\rho,w)$ at every prolongation, so
HAW17 gives $\mathrm{TU}(\pi,w)$ directly.  In both cases
the already chosen $t$ and HAW6 give
$\mathrm{WinAt}(\pi,t,w)$.  This proves the exact quantifier
order and all the clauses of (HAW.31).

The argument retains the original rank, coefficient ring,
auxiliary place, level, point, representation, and ideal.
All odd primes allowed by the original Hida setting,
including $\ell=3$, remain allowed.  The branch $J=\top$
is vacuous at points and is included.  No additional
continuity, good-point, or $\mathrm{V0ok}$ hypothesis has
been added to (HAW.30).
\end{proof}
% END HAW19

\subsection{What this completion supplies}

HAW19 is the entire original Hida $\mathrm{end2}$ field.
HAW10 independently supplies the literal $\mathrm{CoreOrdAlt}$
needed by the reverse point and twist arguments, and HAW13
supplies its full ordered-flag consequence.  The proof retains
the original published foundations listed at the start and
the earlier chapters' named foundation boundaries.  In
particular it does not count a new proof of definite unitary
base change, polarized cuspidal Galois attachment, local
classification, or strong multiplicity one.  The new work
is the explicit human reconstruction of all the finite
comparisons, ordinary extraction, flag and conjugate-place
transport, and strict uniform-window consumers leading to
the literal endpoint.

